\subsection{Magnus 群}

对一切 \(a = (a_n)_{n \geq 0}\) （属于 \(\hat{\mathbf{A}}(\mathbf{X})\) ），元素 \(a_0\) （属于 \(\mathbf{K}\) ）称为 \(a\) 的常数项，记作 \(\varepsilon(a)\) 。公式 (1) 表明 \(\varepsilon\) 是 \(\hat{\mathbf{A}}(\mathbf{X})\) 到 \(\mathbf{K}\) 中的一个代数同态。

引理 1. 元素 \(a\) （属于 \(\hat{\mathbf{A}}(\mathbf{X})\) ）可逆，当且仅当它的常数项在 \(\mathbf{K}\) 中可逆。

若 \(a\) 在 \(\hat{\mathbf{A}}(\mathbf{X})\) 中可逆，则 \(\varepsilon(a)\) 在 \(\mathbf{K}\) 中可逆。反之，若 \(\varepsilon(a)\) 在 \(\mathbf{K}\) 中可逆，则存在 \(u \in \hat{\mathbf{A}}_1(\mathbf{X})\) 使得 \(a = \varepsilon(a)(1 - u)\) ；我们记 \(b = \left(\sum_{n \geqslant 0} u^n\right)\varepsilon(a)^{-1}\) 。那么 \(ab = ba = 1\) ，并且 \(a\) 是可逆的。

因此 \(\hat{\mathbf{A}}(\mathbf{X})\) 中常数项为 1 的元素组成的集合是乘法幺半群 \(\hat{\mathbf{A}}(\mathbf{X})\) 的一个子群，称为在 \(\mathbf{X}\) 上（相对于 \(\mathbf{K}\) ）构造的 Magnus 群。在本章中它将记作 \(\Gamma(\mathbf{X})\) 或简称 \(\Gamma\) 。对每个整数 \(n \geqslant 1\) ，用 \(\Gamma_n\) 表示 \(a \in \Gamma\) 组成的集合，其中 \(\omega(a - 1) \geqslant n\) 。由 § 4 no. 5 的命题 2，序列 \((\Gamma_n)_{n \geqslant 1}\) 是 \(\Gamma\) 上的一个整中心滤过。

